
\documentclass[11pt,a4paper]{article}

% Cargar configuración desde archivo externo
% =========================
% PAQUETES
% =========================
\usepackage[utf8]{inputenc}
\usepackage[T1]{fontenc}
\usepackage[spanish]{babel}
\usepackage{geometry}
\usepackage{graphicx}
\usepackage{xcolor}
\usepackage{tikz}
\usepackage{tabularx}
\usepackage{enumitem}
\usepackage{hyperref}
\usepackage{fancyhdr}
\usepackage{parskip}
\usepackage{ragged2e}
\usepackage{helvet}
\usepackage{titlesec}  % <-- Agregar este paquete
\renewcommand{\familydefault}{\sfdefault}

% Justificacion uniforme y mejor ajuste de lineas para evitar desbordes.
\AtBeginDocument{\justifying}
\setlength{\emergencystretch}{1.5em}

% =========================
% GEOMETRÍA
% =========================
\geometry{
    top=2cm,
    bottom=2.5cm,
    left=2.5cm,
    right=2.5cm
}

% =========================
% COLORES PERSONALIZADOS
% =========================
\definecolor{ieeeOrange}{HTML}{E07830}
\definecolor{ieeeGreen}{HTML}{8FBC8F}
\definecolor{ieeeLightBlue}{HTML}{87CEEB}
\definecolor{ieeeGray}{HTML}{C0C0C0}
\definecolor{titleBrown}{HTML}{5D4E37}
\definecolor{sectionTeal}{HTML}{008080}
\definecolor{linkTeal}{HTML}{007070}
\definecolor{textBrown}{HTML}{716650}

% =========================
% FORMATEO DE SECCIONES
% =========================
\titleformat{\section}
    {\titlerule
    \Large\bfseries\color{ieeeOrange}}  % Formato: grande, negrita, naranja
    {}                                    % Sin número
    {0pt}                                 % Sin espacio entre número y título
    {}                                    % Sin prefijo
\titlespacing*{\section}{0pt}{0.8cm}{0.4cm}  % Espaciado: izq, antes, después

\titleformat{\subsection}
    {\large\bfseries\color{ieeeOrange}}
    {}
    {0pt}
    {}
\titlespacing*{\subsection}{0pt}{0.6cm}{0.3cm}

% =========================
% COMANDOS PERSONALIZADOS
% =========================

% Barra de colores superior
\newcommand{\colorbar}{%
    \noindent
    \begin{tikzpicture}[remember picture, overlay]
        \fill[ieeeOrange] (current page.north west) rectangle ++(5.25cm,-0.4cm);
        \fill[ieeeGreen] ([xshift=5.25cm]current page.north west) rectangle ++(5.25cm,-0.4cm);
        \fill[ieeeLightBlue] ([xshift=10.5cm]current page.north west) rectangle ++(5.25cm,-0.4cm);
        \fill[ieeeGray] ([xshift=15.75cm]current page.north west) rectangle ++(5.25cm,-0.4cm);
    \end{tikzpicture}
}

% Línea de acción (color naranja)
\newcommand{\lineaAccion}[2]{%
    \subsection*{Línea de Acción #1: #2}
}

% Fila de tabla descriptiva
\newcommand{\filaDescriptiva}[2]{%
    \textbf{#1} & #2 \\[0.4cm]
}

\setlist[enumerate,1]{label=\Roman*., font=\bfseries, leftmargin=1.5cm}
\setlist[enumerate,2]{label=\roman*., leftmargin=1cm}

\newcommand{\logos}{
% Logos (placeholder - reemplazar con logos reales)
\noindent
\begin{minipage}[c]{0.4\textwidth}
    \centering
    % Logo IEEE (diamante)
    \includegraphics[width=0.9\linewidth]{src/IEEEBLK.png}
\end{minipage}%
\hfill
\hfill
\begin{minipage}[c]{0.3\textwidth}
    \centering
    % Logo Young Professionals
    \includegraphics[width=1.2\linewidth]{src/UAA2.png}

\end{minipage}



}

% =========================
% DOCUMENTO
% =========================
\begin{document}
\color{textBrown}

% ============================================================
% PORTADA
% ============================================================
\colorbar
\pagestyle{empty}

\vspace{1.5cm}
\logos
\vspace{3cm}

% Título principal
{\fontsize{36}{42}\selectfont\bfseries\color{titleBrown} Protocolo:\\ Atención a Ponentes}

\vspace{0.8cm}

% Subtítulo
{\large Rama Estudiantil IEEE Campus UAA}

\vspace{0.3cm}

% Línea decorativa
\noindent\rule{7.5cm}{0.5pt}

\vfill

% Información inferior
{\color{sectionTeal}\bfseries IEEE Rama UAA}\\[0.2cm]
Última fecha de revisión: Julio 2026\\[0.1cm]

\newpage

% ============================================================
% DATOS DE IDENTIFICACIÓN
% ============================================================
\pagestyle{fancy}
\fancyhf{}
\renewcommand{\headrulewidth}{0pt}
\fancyfoot[C]{\thepage}

\vspace{0.5cm}

\section{Datos de identificación del documento}

\vspace{0.8cm}

\hyphenpenalty=10000
\exhyphenpenalty=10000

\begin{center}
\begin{tabular}{@{}p{4.5cm} p{8cm}@{}}
    \small Categoría & \small Detalle \\[0.5cm]
    
    \textbf{Título del Documento} & Protocolo Integral de Atención a Conferencistas \\[0.5cm]
    
    \textbf{Sección IEEE} & Aguascalientes (Región 9) \\[0.5cm]
    
    \textbf{Coordinación Responsable} & Rama Estudiantil IEEE Campus UAA \\[0.5cm]
    
    \textbf{Liderazgo} & Das Reyes (Vicepresidente) \\[0.5cm]
    
    \textbf{Fecha de Redacción} & Julio de 2026 \\[0.5cm]
    
\end{tabular}
\end{center}

% ============================================================
% CONTENIDO
% ============================================================

\section{Introducción}
El presente protocolo establece los lineamientos estandarizados para la gestión, trato y seguimiento de los conferencistas y ponentes invitados por la Rama Estudiantil IEEE Campus UAA. El objetivo es garantizar una experiencia profesional, cálida y fluida, proyectando la excelencia de nuestra comunidad universitaria y asegurando el éxito de eventos como el \textit{Engineering Showcase} o ciclos de conferencias.
\newpage
\section{Fase 1: Previo al Evento (Contacto y Planeación)}

El trato al ponente comienza desde el primer momento en que se le considera para el evento. Es vital proyectar organización y formalidad.

\subsection{1.1 Asignación de un Enlace (Anfitrión)}
Se debe designar a un único miembro de la mesa directiva (o voluntario capacitado) como el \textbf{Punto de Contacto (PoC)} exclusivo del ponente. Esto evita cruce de correos y confusiones.

\subsection{1.2 El Correo de Primer Contacto}
Debe ser claro, formal y persuasivo. A continuación, se presenta una plantilla base:

\vspace{0.5cm}
\begin{center}
\fbox{
\begin{minipage}{0.9\textwidth}
\vspace{0.3cm}
\textbf{Asunto:} Invitación como Ponente - [Nombre del Evento, ej. Engineering Showcase 2026] \\[0.3cm]
Estimado/a [Grado Académico y Nombre, ej. Ing. Omar Bazán / Dr. Francisco Rangel],\\[0.3cm]
Espero que se encuentre muy bien. Me dirijo a usted en nombre de la \textbf{Rama Estudiantil IEEE de la Universidad Autónoma de Aguascalientes (UAA)}.\\[0.3cm]
Conocemos su destacada trayectoria en [área de especialidad, ej. innovación y desarrollo tecnológico] y el impacto de su trabajo. Por ello, sería un honor para nosotros invitarle a participar como conferencista en nuestro próximo evento, \textbf{[Nombre del Evento]}, el cual se llevará a cabo el día [Fecha] en [Lugar, ej. Auditorio Pedro de Alba].\\[0.3cm]
El objetivo de este espacio es [propósito de la charla, ej. inspirar a los futuros ingenieros]. Consideramos que su visión aportará un valor incalculable a nuestra comunidad estudiantil.\\[0.3cm]
\textbf{Detalles preliminares:}
\begin{itemize}[leftmargin=1cm]
    \item \textbf{Formato:} Presencial, duración de [X] minutos (incluyendo sesión de preguntas y respuestas).
    \item \textbf{Audiencia:} Estudiantes de ingenierías y ciencias básicas.
\end{itemize}
Nos gustaría saber si cuenta con disponibilidad en su agenda. De ser así, podemos agendar una breve llamada para afinar los detalles de su participación.\\[0.3cm]
Agradecemos de antemano su tiempo y consideración.\\[0.3cm]
Atentamente,\\
\textbf{Orlando Reyes}\\
Vicepresidente - Rama Estudiantil IEEE UAA
\vspace{0.3cm}
\end{minipage}
}
\end{center}
\vspace{0.5cm}

\subsection{1.3 Seguimiento y Logística Previa}
Una vez que el ponente acepta, el Enlace debe solicitar y coordinar lo siguiente (al menos 2 semanas antes):
\begin{itemize}
    \item \textbf{Semblanza breve y fotografía:} Para el diseño de flyers y material de difusión.
    \item \textbf{Requerimientos técnicos:} ¿Necesita software específico, adaptadores (HDMI/Type-C), salida de audio, o material extra?
    \item \textbf{Presentación (Diapositivas):} Solicitarla 48 horas antes para tenerla cargada en el equipo base como respaldo ante fallas técnicas.
    \item \textbf{Logística de llegada:} Enviar un mapa del campus central de la UAA, indicaciones de estacionamiento y el número celular del anfitrión. Si el ponente es foráneo, coordinar viáticos y transporte.
    \item \textbf{Recordatorio:} Enviar un mensaje corto de confirmación 24-48 horas antes del evento.
\end{itemize}

\section{Fase 2: Durante el Evento (Recepción y Ejecución)}

La experiencia presencial debe ser impecable para que el ponente se sienta bienvenido y valorado.

\subsection{2.1 Recepción en el Campus}
\begin{itemize}
    \item El Enlace asignado debe esperar al ponente en el punto de acceso acordado (ej. estacionamiento o entrada del edificio) al menos \textbf{30 minutos antes} de su intervención.
    \item Saludar cordialmente, presentarse e identificar al ponente entregándole su gafete de invitado especial.
    \item Guiar al ponente a una sala de espera o zona reservada para evitar que espere de pie. Ofrecer agua, café o algún aperitivo ligero antes de su participación.
\end{itemize}

\subsection{2.2 Pruebas Técnicas}
\begin{itemize}
    \item Realizar prueba de micrófono, apuntador láser y asegurar que la presentación se proyecte correctamente con la relación de aspecto adecuada.
\end{itemize}

\subsection{2.3 Presentación del Ponente al Público}
La persona que actúa como Maestro de Ceremonias debe seguir un guion dinámico. \textbf{Bajo ninguna circunstancia se debe leer el currículum completo}, sino una versión adaptada para conectar con la audiencia.
\begin{itemize}
    \item \textbf{Estructura ideal:} Saludo + Breve descripción del ponente + Trayectoria destacada + Conexión directa del ponente con el tema a tratar + Bienvenida entusiasta.
    \item \textit{Ejemplo:} ``Para hablar sobre los retos de la electrónica actual, hoy nos acompaña un experto que ha liderado proyectos de innovación a nivel nacional. Egresado de... y actual directivo en... Es un honor recibir en la Universidad Autónoma de Aguascalientes al [Nombre]. ¡Recibámoslo con un fuerte aplauso!''
\end{itemize}

\subsection{2.4 Asistencia durante la ponencia}
\begin{itemize}
    \item El staff debe mantenerse visible (generalmente en la primera fila).
    \item Acordar previamente señales visuales para indicar el tiempo restante (ej. letreros de 10 min, 5 min y Tiempo).
    \item Asegurarse de colocar una botella de agua fresca (sin etiqueta, por cuestiones visuales) en el pódium.
\end{itemize}

\section{Fase 3: Post Ceremonia (Cierre y Agradecimiento)}

El cierre marca la impresión final que el conferencista se llevará de la Rama Estudiantil.

\subsection{3.1 Sesión de Preguntas y Respuestas (Q\&A)}
\begin{itemize}
    \item El moderador de IEEE debe retomar el micrófono y coordinar las participaciones. Es altamente recomendable tener una o dos preguntas preparadas por el staff para romper el hielo si el público es tímido.
\end{itemize}

\subsection{3.2 Agradecimiento Institucional y Obsequio}
Al finalizar la sesión de preguntas:
\begin{itemize}
    \item El Maestro de Ceremonias o un representante de la Mesa Directiva sube al escenario.
    \item Se lee en voz alta el reconocimiento impreso. (\textit{``La Rama Estudiantil IEEE de la Universidad Autónoma de Aguascalientes otorga el presente reconocimiento a...''}).
    \item \textbf{Entrega del Reconocimiento y Obsequio:} Entregar el diploma enmarcado o en carpeta institucional, entregando a la par \textbf{un artículo perteneciente a la UAA (como un termo oficial, libreta o merch de la rama)} para fomentar el sentido de pertenencia y que el ponente se lleve un gran recuerdo de la institución.
    \item Tomar la fotografía oficial en el escenario con el público de fondo.
\end{itemize}

\subsection{3.3 Despedida y Seguimiento}
\begin{itemize}
    \item El Enlace acompaña al ponente de regreso a su vehículo, agradeciéndole personalmente.
    \item \textbf{Seguimiento (24 a 48 horas después):} Enviar un correo electrónico formal de agradecimiento. Se sugiere incluir en este correo un enlace a una carpeta con las mejores fotografías del evento para que el ponente pueda compartirlas en sus redes sociales.
\end{itemize}

\end{document}
